\documentclass[10pt,a4paper]{amsart}
\usepackage[margin=19mm]{geometry}
\usepackage{amsmath,amssymb,mathtools}
\usepackage[scr=rsfs,cal=euler]{mathalfa}
\usepackage{xfrac}
\usepackage[unicode,hidelinks,bookmarks=false]{hyperref}
\newcommand{\ie}{i.e.\ }
\newcommand{\cf}{cf.\ }
\newcommand{\resp}{respectively\ }
\newcommand{\Subsection}{\subsection}
\providecommand{\nobreakdash}{\nobreak\hbox{-}}
\setlength{\emergencystretch}{2em}
\multlinegap0pt
\usepackage{bm}
\usepackage{bbm}
\usepackage{dsfont}
\usepackage{xspace}
\usepackage{cancel}
\usepackage{mathtools}
\usepackage{amsthm}
\makeatletter
\@ifundefined{definition}{\newtheorem{definition}{Definition}}{}
\@ifundefined{proposition}{\newtheorem{proposition}[definition]{Proposition}}{}
\@ifundefined{theorem}{\newtheorem{theorem}[definition]{Theorem}}{}
\makeatother
\title[Cohn--Vossen-Type Inequalities for Three-Manifolds and Local]{Cohn--Vossen-Type Inequalities for Three-Manifolds and Locally Conformally Flat Manifolds}
\author{Jialong Deng}
\date{Source: arXiv:2606.01368v2 (CC BY 4.0)}
\begin{document}
\maketitle

\noindent\emph{Source and licence.} Cohn--Vossen-Type Inequalities for Three-Manifolds and Locally Conformally Flat Manifolds. Version arXiv:2606.01368v2 (CC BY 4.0), distributed under the Creative Commons Attribution 4.0 International licence, as recorded in the arXiv licence metadata for this exact version and shown on the abstract page https://arxiv.org/abs/2606.01368v2. Full rights notice: licenses/arxiv-2606.01368-CC-BY-4.0.txt.\par\smallskip

\noindent\textit{Editorial scope.} Counted unit: the Proposition whose source label is \texttt{prop:sign-needed-beta-positive} in the article (source0.tex lines 2894--2905, source label \texttt{prop:sign-needed-beta-positive}), delivered together with the article's own complete original proof of it (source0.tex lines 2907--3075, one \texttt{proof} environment of 169 lines). No mathematics is written, repaired or re-derived: statement and proof are the article's own text line for line. Required context retained uncounted: eq:base-assumptions (lines 2489--2494); thm:infinite-N (lines 2720--2743). Every definition, display and label of those lines is retained unchanged. Omitted from this package and not redistributed: 1--2488, 2495--2719, 2744--2893, 2906--2906, 3076--3093. Nothing inside a retained excerpt is omitted or altered: every retained body is delivered exactly as the article's lines are written, so no omission receipt of any kind accompanies this package. Citations: the retained excerpts carry no citation command at all, so no bibliography entry of the article is transcribed and no reference is redistributed. Reference adaptation: none was needed and none was made; every reference inside the delivered unit resolves inside it, so no unresolved reference or citation is produced. Printed slips kept literal: no printed word, symbol, hypothesis or number of the counted statement or of its proof was altered. Layout and class: the article's own class is \texttt{article} with packages including amsmath, amsthm, amssymb, hyperref, cleveref, geometry, inputenc, esint, stmaryrd, microtype; the wrapper is the lane's pinned \texttt{amsart} editorial wrapper at 10pt on A4, which restates the theorem-like environments the retained text needs and adds the article's own macro definitions that the retained text needs (none), with extra wrapper packages (bm, bbm, dsfont, xspace, cancel, mathtools). The delivered numbering is the unit's own. Assets: no retained span contains a figure environment, a graphics-inclusion command, a PDF-page inclusion, a table or a TikZ picture; no image file is copied into this package, the package declares no asset dependency and no image file is redistributed with it. Licence: version arXiv:2606.01368v2 of this article is distributed under the Creative Commons Attribution 4.0 International licence, as recorded in the arXiv licence metadata for this exact version and shown on the abstract page https://arxiv.org/abs/2606.01368v2. A full rights notice is delivered at licenses/arxiv-2606.01368-CC-BY-4.0.txt. Characters: every retained excerpt is pure ASCII, so the delivered engine, XeLaTeX with its default Latin Modern text font, reports no missing character.
\begin{equation}\label{eq:base-assumptions}
  \mathrm{Ric}_g\geq 0,
  \qquad
  A_p:=\limsup_{r\to\infty}
  r^{2-n}\int_{B(p,r)}\mathrm{Sc}_g\,dV_g<\infty.
\end{equation}

\begin{theorem}[Infinite-dimensional Bakry--\'Emery tensor]
\label{thm:infinite-N}
Let $(M^n,g,e^{-f}\,dV_g)$ ($n\geq 3$)  be a weighted Riemannian manifold
satisfying $\operatorname{Ric}_{f,\infty}\ge 0$, and suppose $g$
satisfies~\eqref{eq:base-assumptions}.
\begin{enumerate}
\item[\textup{(a)}] If $\alpha\le 0$ and $\beta\ge 0$, then for every
$\varepsilon>0$ there exists $r_0=r_0(\varepsilon)$ such that for all
$r\ge r_0$,
\begin{equation}\label{eq:infty-signed-upper}
  \int_{B_g(p,r)}\mathrm{Sc}_{\alpha,\beta}\,dV_g
  \le(1-\alpha)\int_{B_g(p,r)}\mathrm{Sc}_g\,dV_g
  \le(1-\alpha)(A_p+\varepsilon)\,r^{n-2}.
\end{equation}

\item[\textup{(b)}] If $\beta>0$ and $\mathrm{Sc}_{\alpha,\beta}\ge 0$
pointwise, with $\alpha\in\mathbb{R}$ arbitrary, then for every $\varepsilon>0$
there exists $r_0=r_0(\varepsilon)$ such that for all $r\ge r_0$,
\begin{equation}\label{eq:infty-positive-beta}
  \int_{B_g(p,r)}\mathrm{Sc}_{\alpha,\beta}\,dV_g
  \le C(n,\alpha,\beta)(A_p+\varepsilon+1)\,r^{n-2}.
\end{equation}
\end{enumerate}
\end{theorem}

\begin{proposition}[Necessity of the sign condition when $\beta>0$]
\label{prop:sign-needed-beta-positive}
Let $n>2$, $\alpha>0$, and $\beta>0$. The pointwise hypothesis
$\mathrm{Sc}_{\alpha,\beta}\ge 0$ in Theorem~\ref{thm:infinite-N}\textup{(b)}
cannot be omitted: there exists a smooth convex function $f$ on
$(\mathbb{R}^n,g_{E})$ with $\mathrm{Ric}_{f,\infty}\ge 0$
for which $\mathrm{Sc}_{\alpha,\beta}$ changes sign and
\[
  \limsup_{r\to\infty}r^{2-n}
  \int_{B_r(0)}\bigl(\alpha\,\Delta f-\beta|\nabla f|^2\bigr)\,dx=+\infty.
\]
\end{proposition}

\begin{proof}
We work on \((\mathbb R^n,g_{E})\) with base point \(0\). Then
\(\mathrm{Ric}_g=0\), \(\mathrm{Sc}_g=0\), and the standing assumptions hold
with \(A_p=0\). Thus
\[
  \mathrm{Sc}_{\alpha,\beta}
  =
  \alpha\Delta f-\beta|\nabla f|^2 .
\]

We first record the radial identities. Suppose that
\(v\in C^\infty([0,\infty))\), \(v\ge0\), \(v'\ge0\), and \(v\equiv0\) near
\(0\). Define
\[
  F(r):=\int_0^r v(s)\,ds,
  \qquad
  f(x):=F(|x|).
\]
Then \(f\) is smooth and convex on \(\mathbb R^n\). Indeed, for \(r=|x|>0\),
\[
  \nabla^2 f
  =
  v'(r)\,dr\otimes dr
  +
  \frac{v(r)}{r}
  \bigl(g_{E}-dr\otimes dr\bigr)\ge0,
\]
and \(f\) is smooth at the origin because \(v\equiv0\) near \(0\). Hence
\[
  \mathrm{Ric}_{f,\infty}
  =
  \nabla^2 f
  \ge0.
\]
Moreover,
\[
  \Delta f=v'(r)+\frac{n-1}{r}v(r),
  \qquad
  |\nabla f|^2=v(r)^2 .
\]
Writing \(\omega_{n-1}=|\mathbb S^{n-1}|\), the divergence theorem gives
\begin{equation}\label{eq:Phi-short-final}
  \Phi(R):=
  R^{2-n}
  \int_{B_R(0)}
  \bigl(\alpha\Delta f-\beta|\nabla f|^2\bigr)\,dx
  =
  \omega_{n-1}
  \left[
    \alpha Rv(R)
    -
    \beta R^{2-n}\int_0^R v(s)^2s^{n-1}\,ds
  \right].
\end{equation}
The first term is a boundary flux depending only on the current value
\(v(R)\), while the second term is the accumulated \(L^2\)-cost of the
gradient. We exploit this by keeping \(v\) at the old height \(a_{j-1}\) until
just before \(R_j\), and then raising it to the new height \(a_j\) in a thin
shell near \(R_j\).

Set \(R_j:=3j\) and \(a_0:=0\). Choose \(a_j>a_{j-1}\) inductively so that
\begin{equation}\label{eq:aj-final-short}
  \frac{\alpha}{2}a_jR_j
  \ge
  \frac{\beta}{n}a_{j-1}^2R_j^2+j,
  \qquad
  a_j>\frac{\alpha(n-1)}{\beta R_j}.
\end{equation}
Both conditions are lower bounds on \(a_j\), so any sufficiently large \(a_j\)
satisfies them. The first inequality makes the new boundary flux dominate the
old plateau cost by the margin \(j\); the second is used only to force
negativity of \(\mathrm{Sc}_{\alpha,\beta}\) at \(|x|=R_j\).

Choose
\[
  0<\delta_j<
  \min\left\{\frac12,\frac{\alpha}{2\beta a_j}\right\}.
\]
The bound \(\delta_j<1/2\) keeps the shells
\([R_j-\delta_j,R_j]\) pairwise disjoint, since consecutive radii differ by
\(3\). The bound \(\delta_j<\alpha/(2\beta a_j)\) controls the new-shell cost.

Let \(\theta\in C^\infty(\mathbb R)\) be nondecreasing, with
\(\theta\equiv0\) on \((-\infty,0]\) and \(\theta\equiv1\) on
\([1,\infty)\), and define
\[
  v(r):=
  \sum_{j=1}^{\infty}
  (a_j-a_{j-1})
  \theta\!\left(\frac{r-(R_j-\delta_j)}{\delta_j}\right).
\]
The sum is locally finite. Hence \(v\in C^\infty([0,\infty))\), \(v\ge0\),
\(v'\ge0\), and \(v\equiv0\) near \(0\). Therefore the radial function
\(f(x)=F(|x|)\), with \(F(r)=\int_0^r v(s)\,ds\), is smooth and convex on
\(\mathbb R^n\).

At \(R_j\), the first \(j\) layers have completed and no later layer has
begun; the sum therefore telescopes. Thus
\[
  v(R_j)=\sum_{k=1}^j(a_k-a_{k-1})=a_j,
  \qquad
  v'(R_j)=0,
  \qquad
  v(s)\le a_{j-1}\quad\text{for }0\le s\le R_j-\delta_j.
\]
Splitting the cost integral at \(R_j-\delta_j\), and using \(v\le a_{j-1}\)
before the shell and \(v\le a_j\) on the shell, gives
\[
  \int_0^{R_j}v(s)^2s^{n-1}\,ds
  \le
  a_{j-1}^2\int_0^{R_j}s^{n-1}\,ds
  +
  a_j^2\delta_jR_j^{n-1}
  =
  \frac{a_{j-1}^2R_j^n}{n}
  +
  a_j^2\delta_jR_j^{n-1}.
\]
Substituting this estimate into \eqref{eq:Phi-short-final} gives
\[
\begin{aligned}
  \Phi(R_j)
  &\ge
  \omega_{n-1}
  \left[
    \alpha a_jR_j
    -
    \frac{\beta}{n}a_{j-1}^2R_j^2
    -
    \beta a_j^2\delta_jR_j
  \right]  \\
  &\ge
  \omega_{n-1}
  \left[
    \frac{\alpha}{2}a_jR_j
    -
    \frac{\beta}{n}a_{j-1}^2R_j^2
  \right]
  \ge
  \omega_{n-1}j.
\end{aligned}
\]
Here the second line uses
$\beta a_j^2\delta_jR_j\le \frac{\alpha}{2}a_jR_j,$
and the last inequality uses \eqref{eq:aj-final-short}. Hence
\[
  \limsup_{R\to\infty}
  R^{2-n}
  \int_{B_R(0)}
  \bigl(\alpha\Delta f-\beta|\nabla f|^2\bigr)\,dx
  =
  +\infty .
\]

Finally, since \(\Phi(R_j)>0\), the continuous function
\(\mathrm{Sc}_{\alpha,\beta}\) is positive somewhere in \(B_{R_j}(0)\). On the
other hand, at \(|x|=R_j\), one has \(v(R_j)=a_j\) and \(v'(R_j)=0\), so
\[
  \mathrm{Sc}_{\alpha,\beta}
  =
  \alpha\left(v'(R_j)+\frac{n-1}{R_j}v(R_j)\right)
  -
  \beta v(R_j)^2
  =
  \frac{\alpha(n-1)a_j}{R_j}-\beta a_j^2<0
\]
by \eqref{eq:aj-final-short}. Therefore
\(\mathrm{Sc}_{\alpha,\beta}\) changes sign.
\end{proof}

\end{document}
